% ============================================================
% CUMCM 国赛 D 题数据类模板 — v7.2
% 编译:xelatex main.tex(跑两遍)
% 金标准:六段子结构 / 三段式公式 / 四重检验 / 算法对比 / 创新量化
% 核心强化:统计模型 / 参数优化 / 拟合精度+MC
% ============================================================
\documentclass[12pt,a4paper]{ctexart}

\usepackage{amsmath,amssymb,amsfonts}
\usepackage{graphicx}
\usepackage{booktabs}
\usepackage{geometry}
\usepackage{float}
\usepackage{caption}
\usepackage{subcaption}
\usepackage{listings}
\usepackage{hyperref}
\usepackage{enumitem}
\usepackage{xcolor}
\usepackage{indentfirst}
\usepackage{titlesec}
\usepackage{fancyhdr}
\usepackage{setspace}

\geometry{top=2.5cm, bottom=2.5cm, left=2.5cm, right=2.5cm, headheight=14pt}
% ============ 排版规范（对齐国赛优秀论文 B159） ============
% 行距 1.5 倍，首行缩进 2 字符
\linespread{1.5}\selectfont
\setlength{\parindent}{2em}

% 章节标题居中，二级标题左对齐
\titleformat{\section}{\centering\large\bfseries\heiti}{\thesection}{0.8em}{}
\titleformat{\subsection}{\normalsize\bfseries\heiti}{\thesubsection}{0.8em}{}
\titleformat{\subsubsection}{\normalsize\bfseries\heiti}{\thesubsubsection}{0.6em}{}

% 页眉空、页脚页码居中
\pagestyle{fancy}
\fancyhf{}
\fancyfoot[C]{\small\thepage}
\renewcommand{\headrulewidth}{0pt}


\setCJKmainfont{SimSun}[BoldFont=SimHei, ItalicFont=KaiTi]
\setCJKsansfont{SimHei}
\setCJKmonofont{FangSong}

% 图/表题自动编号前缀（ctexart 默认不设置，缺失会导致图题无「图」字）
\renewcommand{\figurename}{图}
\renewcommand{\tablename}{表}


\DeclareMathOperator{\erf}{erf}
\DeclareMathOperator{\sgn}{sgn}

\hypersetup{colorlinks=true, linkcolor=black, citecolor=black, urlcolor=blue}

\lstset{
  basicstyle=\ttfamily\small,
  breaklines=true,
  frame=single,
  numbers=left,
  numberstyle=\tiny,
  keywordstyle=\color{blue},
  commentstyle=\color{green!50!black},
}

\title{\textbf{[研究对象]的数据分析与预测建模}}
\author{队伍编号:\textbf{XXXX}}
\date{}

\begin{document}

\maketitle

% ============================================================
\begin{abstract}
\noindent
\textbf{第一段}:行业背景 + 研究对象 + 三层框架(描述-建模-预测)。

\textbf{第二段}:问题一基于 [N 条] 数据做 EDA,建立 [模型],预测 [对象] 在 [时间] 的值为 [数值](95\%CI:[a, b]);问题二...;问题三...

\textbf{第三段}:3 项创新(特征工程 / 模型改进 $R^2$ 提升 $X\%$ / 预测精度提升 $Y\%$)+ 四重检验指标 + 跨行业迁移。

\vspace{0.5em}
\noindent\textbf{关键词}:[对象] [数据] [模型] [预测] [精度]
\end{abstract}

\newpage
\newpage

% ============================================================
\section{问题重述}
\subsection{问题背景}
[背景:聚焦数据来源与预测需求]
\subsection{问题提出}
\begin{itemize}
  \item 问题一:[数据分析与关系挖掘]
  \item 问题二:[预测建模]
  \item 问题三:[参数优化与情景分析]
\end{itemize}

% ============================================================
\section{问题分析}
\begin{figure}[H]
\centering
\includegraphics[width=0.85\textwidth]{../figures/pdf/fig1_eda.pdf}
\caption{图1 数据特征分布与相关性热力图($x_1$ 与 $x_4$ 强相关 $r=0.82$)}
\label{fig:eda}
\end{figure}

% ============================================================
\section{模型假设}
\begin{enumerate}
  \item 假设1:数据采样具有代表性
  \item 假设2:缺失值处理不引入系统性偏差
  \item 假设3:未来趋势遵循历史规律
  \item 假设4:残差服从独立同分布
\end{enumerate}

% ============================================================
\section{符号说明}
\begin{table}[H]
\centering
\caption{符号说明}
\label{tab:symbols}
\begin{tabular}{llll}
\toprule
符号 & 含义 & 取值范围 & 单位 \\
\midrule
$y_t$ & 目标变量在第 $t$ 期观测值 & - & - \\
$\hat{y}_t$ & 预测值 & - & - \\
$\beta_k$ & 第 $k$ 个回归系数 & - & - \\
$\epsilon_t$ & 残差 & $\mathcal{N}(0, \sigma^2)$ & - \\
\bottomrule
\end{tabular}
\end{table}

% ============================================================
\section{模型建立与求解}

\subsection{问题一:数据分析与关系挖掘}

\subsubsection{参数配置与数据预处理}
[数据清洗:缺失值处理、异常值检测、归一化]

\subsubsection{基础经典模型:线性回归}
为挖掘 [变量] 间关系,建立多元线性回归模型:
\begin{equation}
y_t = \beta_0 + \beta_1 x_{1,t} + \beta_2 x_{2,t} + \cdots + \beta_k x_{k,t} + \epsilon_t
\label{eq:linear_reg}
\end{equation}
其中 $\beta_k$ 为回归系数,$\epsilon_t \sim \mathcal{N}(0, \sigma^2)$ 为残差。式~\eqref{eq:linear_reg} 假设变量间线性关系,需后续检验。

\subsubsection{传统模型缺陷剖析}
\begin{enumerate}
  \item 缺陷1:线性假设过强,$R^2$ 仅 $0.72$
  \item 缺陷2:未捕捉时序依赖,残差自相关 DW=1.2
  \item 缺陷3:对非线性关系建模能力弱
\end{enumerate}

\subsubsection{创新改进模型:ARIMA-LSTM 组合}
本文建立 ARIMA-LSTM 组合预测模型,分别建模线性与非线性成分:
\paragraph{ARIMA 线性成分}
\begin{equation}
\Phi(B) (1-B)^d y_t = \Theta(B) \epsilon_t
\label{eq:arima}
\end{equation}
其中 $B$ 为后移算子,$\Phi, \Theta$ 为 AR/MA 多项式。

式~\eqref{eq:arima} 建模时序线性成分,其中 $B$ 为后移算子。

\paragraph{LSTM 非线性残差}
\begin{equation}
\hat{r}_{t+1} = \text{LSTM}(r_{t-L+1}, \ldots, r_t; \mathbf{W}, \mathbf{b})
\label{eq:lstm}
\end{equation}
式~\eqref{eq:lstm} 用 LSTM 建模非线性残差,其中 $r_t = y_t - \hat{y}_{\text{ARIMA},t}$ 为 ARIMA 残差。组合预测 $\hat{y}_t = \hat{y}_{\text{ARIMA},t} + \hat{r}_t$。

\subsubsection{算法设计与求解}

\paragraph{三算法对比表}
\begin{table}[H]
\centering
\caption{预测算法对比}
\label{tab:algo_compare}
\begin{tabular}{lccccc}
\toprule
算法 & $R^2$ & RMSE & MAPE & 训练时间 (s) & 评级 \\
\midrule
ARIMA-LSTM & 0.95 & 0.18 & 2.1\% & 320 & ★★★ \\
纯 LSTM & 0.89 & 0.27 & 3.5\% & 280 & ★★☆ \\
纯 ARIMA & 0.78 & 0.42 & 5.8\% & 12 & ★☆☆ \\
\bottomrule
\end{tabular}
\end{table}

选用 ARIMA-LSTM 组合的理由:$R^2 = 0.95$ 最高,RMSE $= 0.18$ 最低,兼顾线性和非线性建模能力。

\subsubsection{数值结果与可视化}

\begin{figure}[H]
\centering
\includegraphics[width=0.85\textwidth]{../figures/pdf/fig2_prediction.pdf}
\caption{图2 预测值与真实值对比(95\%CI 阴影带, $R^2 = 0.95$)}
\label{fig:prediction}
\end{figure}

\begin{table}[H]
\centering
\caption{问题一预测结果}
\label{tab:result1}
\begin{tabular}{lcccc}
\toprule
时间 & 真实值 & 预测值 & 95\%CI & 相对误差 \\
\midrule
$t+1$ & 52.3 & 52.1 & [51.4, 52.8] & 0.4\% \\
$t+2$ & 54.8 & 54.5 & [53.6, 55.4] & 0.5\% \\
$t+3$ & 56.2 & 56.0 & [54.8, 57.2] & 0.4\% \\
\bottomrule
\end{tabular}
\end{table}

图~\ref{fig:prediction} 显示组合模型预测值紧贴真实值,95\%CI 覆盖率 $96.7\%$,较纯 ARIMA 的 $R^2$ 提升 $0.17$。

\subsection{问题二}
% 同六段子结构

\subsection{问题三}
% 同六段子结构

% ============================================================
\section{模型检验}

\subsection{拟合精度检验}
$R^2 = 0.95$, $\text{MAE} = 0.15$, $\text{RMSE} = 0.18$, $\text{MAPE} = 2.1\%$。残差 Shapiro-Wilk $p = 0.78$,Ljung-Box $p = 0.65$(无自相关)。

\begin{figure}[H]
\centering
\includegraphics[width=0.85\textwidth]{../figures/pdf/figN_residual.pdf}
\caption{图N 残差诊断(残差符合正态分布且无自相关)}
\label{fig:residual}
\end{figure}

\subsection{灵敏度分析}
\begin{table}[H]
\centering
\caption{模型参数灵敏度分析}
\label{tab:sensitivity}
\begin{tabular}{lccccc}
\toprule
参数 & 基准值 & ±10\% 影响 & ±20\% 影响 & 弹性系数 $\varepsilon$ & 分级 \\
\midrule
学习率 $\eta$ & 0.001 & ±0.8\% & ±1.5\% & 0.08 & 低 \\
窗口 $L$ & 12 & ±2.3\% & ±4.8\% & 0.24 & 中 \\
AR 阶数 $p$ & 3 & ±3.5\% & ±7.2\% & 0.36 & 中 \\
\bottomrule
\end{tabular}
\end{table}

\begin{figure}[H]
\centering
\includegraphics[width=0.85\textwidth]{../figures/pdf/figN_tornado.pdf}
\caption{图N 灵敏度龙卷风图(AR 阶数 $p$ 为中灵敏度参数)}
\label{fig:tornado}
\end{figure}

\subsection{蒙特卡洛验证}
$N_{\text{sim}} = 200$ 次随机扰动训练数据,预测均值 $\mu = 52.1$,标准差 $\sigma = 0.8$,$\text{CV} = 1.5\%$,$95\%\text{CI} = [50.5, 53.7]$。

\begin{figure}[H]
\centering
\includegraphics[width=0.85\textwidth]{../figures/pdf/figN_mc.pdf}
\caption{图N 蒙特卡洛仿真分布(200 次, $\mu=52.1$, 95\%CI=[50.5, 53.7])}
\label{fig:mc}
\end{figure}

\subsection{假设误差量化}
\begin{table}[H]
\centering
\caption{假设误差量化}
\label{tab:assumption_error}
\begin{tabular}{lcc}
\toprule
假设 & 偏差量化 & 对预测影响 \\
\midrule
假设1:数据代表性 & 分布偏度 $\gamma = 0.32$ & 预测偏差 $1.2\%$ \\
假设2:缺失值无偏 & 缺失率 $3.5\%$ & 预测偏差 $0.8\%$ \\
假设3:趋势延续 & 趋势变动 $\Delta = 5\%$ & 预测偏差 $2.1\%$ \\
假设4:残差独立 & DW = 1.95 & 接近理想 \\
\bottomrule
\end{tabular}
\end{table}

% ============================================================
\section{模型优缺点与改进}
\textbf{创新点}:
\begin{enumerate}
  \item 构建特征工程方案,提取 [N 个] 衍生特征,$R^2$ 提升 $0.17$
  \item ARIMA-LSTM 组合模型,较纯 LSTM 精度提升 $0.06$,较纯 ARIMA 提升 $0.17$
  \item 蒙特卡洛置信区间量化预测不确定性,可靠性 $96.7\%$
\end{enumerate}

% ============================================================
\section{参考文献}
\begin{thebibliography}{99}
\bibitem{ref1} 作者. 标题[J]. 期刊, 年份, 卷(期): 页码.
\bibitem{ref2} Author A. Title[J]. Journal, Year, Vol: Pages.
% 近5年 ≥40%, 外文 ≥30%, 总数 ≥10
\end{thebibliography}

% ============================================================
\appendix
\section{附录1 数据说明}
[数据字段说明、缺失值处理、result 样表]

\section{附录2 代码清单}
% ============ 国一铁律：附录必须含完整代码，缺代码取消评奖资格 ============
% 完整代码用 \lstinputlisting 从 code/ 目录读入，禁止只贴片段/伪代码。
% 每个脚本附输入/输出/功能说明，确保可复现（全部固定 SEED=42）。
\begin{table}[H]
\centering
\caption{代码模块说明（完整代码见附件，固定 SEED=42 可复现）}
\label{tab:code_modules}
\begin{tabular}{llll}
\toprule
脚本 & 功能 & 输入 & 输出 \\
\midrule
\texttt{code/00\_data\_prep.py} & 数据预处理与质量校验 & 附件1.xlsx & \texttt{data/clean/*.csv} \\
\texttt{code/01\_q1\_analysis.py} & 问题一求解 & 清洗后数据 & \texttt{results/q1\_*.json} \\
\texttt{code/02\_q2\_classify.py} & 问题二求解 & 关键词表 & \texttt{results/q2\_*.json} \\
\texttt{code/03\_q3\_optimize.py} & 问题三求解 & 问题二结果 & \texttt{results/q3\_*.json} \\
\texttt{code/04\_q4\_robust.py} & 问题四求解 & 问题三模型 & \texttt{results/q4\_*.json} \\
\texttt{code/05\_figures.py} & 图表生成 & \texttt{results/*.json} & \texttt{figures/png}, \texttt{figures/pdf} \\
\bottomrule
\end{tabular}
\end{table}

% 逐文件引入完整代码（示例）：
% \lstinputlisting[language=Python, caption=code/01_q1_analysis.py]{../code/01_q1_analysis.py}


\section{附录3 公式推导}
[ARIMA-LSTM 组合模型完整推导]

\end{document}
